\documentclass[11pt,a4paper]{article}

\usepackage[utf8]{inputenc}
\usepackage{amsmath, amssymb, amsfonts}
\usepackage{graphicx}
\usepackage{hyperref}
\usepackage{geometry}
\usepackage{listings}
\usepackage{xcolor}
\usepackage{cite}
\usepackage{booktabs}

\geometry{a4paper, margin=1in}

% Code listing style for Julia
\lstdefinelanguage{Julia}%
  {morekeywords={abstract,break,case,catch,const,continue,do,else,elseif,%
      end,export,false,for,function,immutable,import,importall,if,in,%
      macro,module,otherwise,quote,return,switch,true,try,type,typealias,%
      using,while},%
   sensitive=true,%
   alsoother={$},%
   morecomment=[l]\#,%
   morecomment=[n]{\#=}{=\#},%
   morestring=[s]{"}{"},%
   morestring=[m]{'}{'},%
}[keywords,comments,strings]%

\lstset{%
    language         = Julia,
    basicstyle       = \ttfamily\small,
    keywordstyle     = \bfseries\color{blue},
    stringstyle      = \color{magenta},
    commentstyle     = \color{green!50!black},
    showstringspaces = false,
    numbers          = none,
    numberstyle      = \tiny\color{gray},
    stepnumber       = 1,
    numbersep        = 10pt,
    tabsize          = 4,
    showspaces       = false,
    showtabs         = false,
    frame            = single,
    rulecolor        = \color{black!30},
    backgroundcolor  = \color{gray!5},
    breaklines       = true
}

\title{\textbf{Learning High-Fidelity Monte Carlo Dosimetry from SPECT/CT using Universal Differential Equations for $^{177}$Lu-PSMA Therapy}}
\author{Scientific Machine Learning (SciML) Implementation Guide}
\date{\today}

\begin{document}

\maketitle

\begin{abstract}
Targeted radiopharmaceutical therapy (TRT) using $^{177}$Lu-PSMA requires highly accurate, patient-specific voxel-level dosimetry. This paper presents a exhaustive programmatic framework using \texttt{DifferentialEquations.jl} and the SciML ecosystem to formulate a Universal Differential Equation (UDE). By embedding a neural network corrector within the radiation transport equations, we rapidly learn the complex transformation from homogeneous kernel convolutions to high-fidelity MC dosemaps. This implementation explicitly incorporates all established physical, biological, and imaging variables for $^{177}$Lu-PSMA, including system-specific calibration, saturable receptor kinetics, and biphasic clearance.
\end{abstract}

\section{Introduction}
Scientific Machine Learning (SciML) offers a path to bridge the gap between fast but inaccurate analytical models and slow gold-standard Monte Carlo (MC) simulations by combining \textbf{Universal Differential Equations (UDEs)} with patient-specific anatomical (CT) and functional (SPECT) data \cite{KennedyLugassi2020}.

\section{Learning Objectives: Known Physics vs. Data-Driven Discovery}

The UDE framework partitions the dosimetry problem into \textit{Mechanistic Knowns} and \textit{Learned Uncertainties}.

\subsection{Compartmental Pharmacokinetics (PK)}
\textbf{What is Known (Mechanistic Constraints):}
\begin{itemize}
    \item \textbf{Radiological Constants:} Physical decay constant $\lambda_{phys} = 4.34 \times 10^{-3} \text{ h}^{-1}$ and mean energy per decay $\bar{E}_i$ \cite{ReschFouladgar2023, StabinMadsen2019}.
    \item \textbf{Conservation of Mass:} Mass-balance equations for tracer movement between the central blood compartment and tissue voxels. For $^{177}$Lu-PSMA, the \textbf{blood circulation rate} $\lambda_{bc}$ is defined as $\ln(2)/1 \text{ min} \approx 41.58 \text{ h}^{-1}$, representing a circulation half-life of 1 minute to model the rapid initial organ uptake post-injection \cite{HardiansyahYousefzadehNowshahr2024}.
    \item \textbf{Renal Clearance Scaling:} Excretion rate $k_{10}$ scales with individual Glomerular Filtration Rate (GFR) \cite{HardiansyahYousefzadehNowshahr2024}.
    \item \textbf{Imaging Corrections:} Initial activity scaling via the \textbf{Camera Calibration Factor} ($CF$) and resolution recovery via \textbf{Recovery Coefficients} ($RC$). $CF$ is derived from DICOM tags \texttt{(0054,1001)} (Units, e.g., Bq/ml or counts/s/MBq) and quantified using \texttt{(0028,1053)} (Rescale Slope) and \texttt{(0028,1052)} (Rescale Intercept) to convert raw counts to activity \cite{GleisnerChouin2022}. $RC$ values are applied as structure-specific look-up tables based on object size (e.g., from NEMA phantom measurements) to correct for Partial Volume Effects (PVE) \cite{KlettingSchimmel2015}.
\end{itemize}

\textbf{What needs to be Learned (Neural Targets):}
\begin{itemize}
    \item \textbf{Voxel-Wise $B_{MAX}$:} Local receptor saturation thresholds defining the \textit{Tumour Sink Effect} \cite{BegumThieme2018}.
    \item \textbf{Kinetic Drift:} Cycle-dependent shifts in uptake rates ($k_{in}, k_{out}$) due to tumour regression \cite{OkamotoThieme2016, SiebingaVeen2024}.
\end{itemize}

\subsection{Radiation Transport and Scattering}
\textbf{What is Known (Mechanistic Constraints):}
\begin{itemize}
    \item \textbf{Base Physics:} \textbf{Energy Deposition Kernels} (Dual $K_{\beta, lung}$ and $K_{\gamma, water}$) and \textbf{Hounsfield Unit (HU) to mass density} ($\rho$) conversion \cite{GtzSchmidkonz2019}.
    \item \textbf{Energy Deposition Kernel ($K_{base}$):} To ensure maximum physical geometry tracking for $1.8 \text{ mm}$ to $6.0 \text{ mm}$ beta traversal, $K_{\beta, lung}$ is a dimensionless 3D probability matrix calculated via Monte Carlo simulations in a lung-equivalent medium ($\rho \approx 0.3 \text{ g/cm}^3$). If purely water gradients are used natively, the base geometry completely truncates any radiation calculation in lower-density lung tissues ($0 \cdot \exp = 0$). Furthermore, to accurately capture distant organ cross-fire and prevent premature truncation of the 113/208 keV gamma emissions (representing roughly 10-15\% of total absorbed dose), a secondary broad-range gamma kernel ($K_{\gamma, water}$) is summed into the convolutional base. By establishing dual initial convolution matrices ($\dot{\text{Energy}}_{base} = \tilde{A} \ast (\bar{E}_{\beta} K_{\beta, lung} + \bar{E}_{\gamma} K_{\gamma, water})$), the "baseline" implicitly expands to track both short-range scalar momentum and deep-tier photon projection geometry natively.
    \item \textbf{HU-to-Density ($\rho$) Conversion:} Anatomical CT data is transformed into a voxel-wise mass density map ($\rho(\mathbf{r})$) using a parameterized \textbf{bi-linear calibration curve}. This curve maps HU values to physical density based on scanner-specific phantom measurements, typically using two linear segments with a breakpoint at the density of water ($HU=0, \rho=1.0$). By default, reasonable constants ($\text{slope}_{air} = 0.001$, $\text{slope}_{tissue} = 0.0007$) are provided, but these should be tailored to specific hospital hardware:
    \begin{equation}
        \rho = \begin{cases} 1.0 + \text{slope}_{air} \cdot HU & HU \leq 0 \text{ (Air to Water)} \\ 1.0 + \text{slope}_{tissue} \cdot HU & HU > 0 \text{ (Water to Bone)} \end{cases}
    \end{equation}
    This density map is essential for calculating the target mass $m(\mathbf{r})$ and serves as a primary spatial input for the neural corrector $\mathcal{N}_\theta$ \cite{GtzSchmidkonz2019, BaileyMarquis2022}.
    \item \textbf{Energy-to-Dose Normalization:} The fundamental definition of absorbed dose is energy deposited per unit mass ($1 \text{ Gy} = 1 \text{ J/kg}$). In the UDE framework, the transport sub-system initially calculates the total energy deposited in a voxel. To convert this to a clinically relevant dose rate $\dot{D}$ (Gy/h), the energy must be normalized by the \textbf{target voxel mass} $m(\mathbf{r})$:
    \begin{equation}
        m(\mathbf{r}) = V_{voxel} \cdot \rho(\mathbf{r})
    \end{equation}
    where $V_{voxel}$ is the constant physical volume of the voxel (derived from DICOM Pixel Spacing and Slice Thickness) and $\rho(\mathbf{r})$ is the spatially varying density map derived from the CT HU conversion. This voxel-specific normalization is critical for accurately modeling energy deposition in heterogeneous tissues, such as the low-mass lung ($ \rho \approx 0.3 \text{ g/cm}^3$) versus high-mass cortical bone ($\rho \approx 1.9 \text{ g/cm}^3$), where identical energy deposition results in vastly different absorbed doses \cite{GtzSchmidkonz2019, StabinMadsen2019}.
\end{itemize}

\textbf{What needs to be Learned (Neural Targets):}
\begin{itemize}
    \item \textbf{Heterogeneous Scattering and Directional Backscatter:} The corrector $\mathcal{N}_\theta$ learns residuals caused by density gradients $\nabla \rho$ and mass-energy absorption variations $\mu_{en}/\rho$ \cite{BaileyMarquis2022, MikellKappadath2016}. To facilitate this without generating directional blindness toward incoming trajectories, the neural network explicitly evaluates four input branches: the scalar baseline energy field ($E_{base}$), the local physical density space ($\rho$), spatial density boundary gradients ($\nabla \rho$) computed via continuous central differences, and crucially, an explicit First-Order Directional Vector Flux field ($\vec{\Phi}_{ray}(\mathbf{r})$). This vector flux field is extracted via unscattered ray-tracing algorithms over the $\tilde{A}$ manifold, enforcing natively required backscattering directions where simple un-vectored gradients blind traditional algorithms. Activity ($A_{total}$) is intentionally factored out to strictly enforce the radiation superposition principle.
\end{itemize}

\section{Radionuclide and Pharmacokinetic Constants for $^{177}$Lu-PSMA}

\subsection{Physical and Radiological Constants}
\begin{itemize}
    \item \textbf{Physical Half-life ($T_{1/2,phys}$):} $\approx 6.647$ days (159.5 hours) \cite{HardiansyahYousefzadehNowshahr2024}.
    \item \textbf{Physical Decay Constant ($\lambda_{phys}$):} $4.34 \times 10^{-3} \text{ h}^{-1}$ \cite{ReschFouladgar2023}.
    \item \textbf{Mean Photon Energies:} Primary peaks at 113 keV and 208 keV \cite{StabinMadsen2019}.
    \item \textbf{Mean Beta ($\beta^-$) Energy:} $\approx 133.6$ keV ($E_{max} = 497$ keV). This is the primary driver of local therapeutic absorbed dose.
    \item \textbf{Absorbed Dose per Activity (Kidneys):} $0.49 - 0.88$ Gy/GBq \cite{ChuamsaamarkkeeCharoenphun2025, OkamotoThieme2016}.
\end{itemize}

\subsection{Organ-Specific Kinetic Constants}
\begin{table}[h!]
\centering
\caption{Kinetic Parameters for $^{177}$Lu-PSMA-617 and $^{177}$Lu-PSMA-I\&T.}
\begin{tabular}{@{}llll@{}}
\toprule
\textbf{Target Region} & \textbf{Effective $T_{1/2}$ (h)} & \textbf{Uptake Rate $k_{in}$ (h$^{-1}$)} & \textbf{Washout Rate $k_{out}$ (h$^{-1}$)} \\ \midrule
Tumours (PSMA-617) & 50 -- 72 \cite{PetersHofferber2021} & 0.000248 \cite{SiebingaPriv2023} & 0.00902 \cite{SiebingaPriv2023} \\
Tumours (PSMA-I\&T) & 30 -- 55 \cite{SiebingaVeen2024} & 0.00967 \cite{SiebingaVeen2024} & 0.0247 \cite{SiebingaVeen2024} \\
Kidneys (PSMA-617) & 28 -- 39 \cite{PetersHofferber2021} & 0.00867 \cite{SiebingaPriv2023} & 0.0141 \cite{SiebingaPriv2023} \\
Salivary Glands & 32.5 -- 42 \cite{PetersPriv2021} & 0.0238 \cite{SiebingaPriv2023} & 0.0140 \cite{SiebingaPriv2023} \\
\bottomrule
\end{tabular}
\end{table}

\section{Exhaustive UDE Model Architecture}

\subsection{Initial Activity Calibration}
The observed SPECT activity $A_{obs}(\mathbf{r}, 0)$ must be scaled before integration. However, modern quantitative SPECT reconstructions natively embed absolute quantification directly into the \texttt{Rescale Slope} (DICOM Tag 0028,1053) yielding Bq/mL. A strict conditional check on the Units Tag \texttt{(0054,1001)} is required to prevent extreme under-dosing via double-calibration. If units are natively \texttt{BQML}, $CF$ strictly defaults to $1.0$:
\begin{equation}
    A_{true}(\mathbf{r}, 0) = \frac{\left( A_{obs}(\mathbf{r}, 0) \cdot \text{Rescale Slope} \right) + \text{Rescale Intercept}}{CF \cdot RC(\mathbf{r})}
\end{equation}
Furthermore, applying an empirical Recovery Coefficient ($RC(\mathbf{r})$) dynamically at the voxel level is statistically unstable and risks violently amplifying high-frequency Gibbs artifacts at lesion edges; instead, it must be restricted strictly to \textit{bounding ROI volumes}.

\subsection{The Govering System of Equations}
\textbf{1. PK Sub-system (Voxel-Wise Surface Binding and Endosomal Internalization):}
\begin{align}
    \frac{dA_{surface}}{dt} &= k_3 A_{free} \left( 1 - \frac{M_{bound}}{B_{MAX}(\mathbf{r})} \right) - (k_4 + k_5 + \lambda_{phys}) A_{surface} \\
    \frac{dA_{internal}}{dt} &= k_5 A_{surface} - \lambda_{phys} A_{internal}
\end{align}
(Where the Molar concentration drives true stoichiometric saturation across both attached domains $M_{bound}(t) = \frac{A_{surface}(t) + A_{internal}(t)}{SA_{MBq/nmol}} \cdot e^{\lambda_{phys} t}$)

\textbf{2. Transport Sub-system (Absorbed Dose Rate and BED Radiobiology):}
Because 3D FFT convolution inside the temporal ODE is mathematically intractable, we decouple the domains. The continuous ODE strictly integrates a macroscopic Time-Integrated Activity loop, explicitly adding a fractional vascular term ($v_b$) to ensure highly perfused marrow and organs receive blood-pool cross-fire during spatial convolutions:
\begin{equation}
    \frac{dTIA}{dt} = A_{free} + A_{surface} + A_{internal} + v_b(\mathbf{r}) \cdot V_{voxel} \cdot \frac{A_{blood}}{V_{blood}}
\end{equation}
The static TIA map ($\tilde{A}$) is then convolved to calculate the physical baseline energy spanning Beta limits and Gamma bounds ($E_{base, total} = \tilde{A} \ast (\bar{E}_\beta K_{\beta, lung} + \bar{E}_\gamma K_{\gamma, water})$). To properly resolve cross-fire from adjacent sources into heterogeneous target masses without breaking trajectory knowledge, the neural network ($\mathcal{N}_\theta$) evaluates the combined baseline energy, physical density features, and explicit $\vec{\Phi}_{ray}$ directional vectors to generate a multiplicative specific absorbed fraction scalar. Note the exponential mapping, ensuring a strict identity map limit $\exp(0)=1.0$ at epoch zero:
\begin{equation}
    D_{phys}(\mathbf{r}) = \frac{1}{m(\mathbf{r})} \cdot \left[ E_{base,total}(\mathbf{r}) \cdot \exp \Big(\mathcal{N}_\theta \big(E_{base,total}, \rho, \nabla \rho, \vec{\Phi}_{ray} \big)\Big) \right]
\end{equation}
To prevent the Neural Corrector from acting as a non-physical energy sink (violating the First Law of Thermodynamics by arbitrarily scaling down physical dose arrays unconstrained), the SciML reverse-mode automatic differentiation incorporates a global spatial conservation limit ($\lambda_{EC}$) constraining the network evaluation exclusively to native geometric re-routing:
\begin{equation}
    \lambda_{EC} \left\| \int D_{phys}(\mathbf{r}) m(\mathbf{r}) dV - \int E_{base}(\mathbf{r}) dV \right\|^2 \to 0
\end{equation}
Because the exact temporal evolution of the dose rate ($\dot{D}(t)$) is lost during spatial decoupling, continuous Markovian Sublethal Damage tracing is abandoned in favor of a post-hoc analytical Biologically Effective Dose (BED) approximation mapping true Sublethal Repair (SLDR) with an absolute macroscopic tumor repopulation subtraction penalty:
\begin{equation}
    \text{BED}(\mathbf{r}) \approx D_{phys}(\mathbf{r}) \cdot \left[ 1 + \frac{D_{phys}(\mathbf{r})}{\alpha/\beta} \cdot \frac{T_{1/2, repair}}{T_{1/2, eff} + T_{1/2, repair}} \right] - \frac{\ln(2)}{\alpha} \frac{T_{treat}(\mathbf{r})}{T_{pot}}
\end{equation}
Because Lu-177 imposes a protracted decay trace, applying a global constant repopulation threshold massively over-penalizes healthy tissues or quickly clearing target sites. Mathematically, the BED formula strictly tracks $T_{treat}(\mathbf{r})$ individually bounded per voxel—defining exactly the absolute duration where exponential local dose rate ($\dot{D}(t, \mathbf{r})$) strictly exceeds biological proliferation index boundaries ($\dot{D}_{lethal}$).

\section{SciML Implementation Details}

\begin{lstlisting}[language=Julia]
function universal_dosimetry_ude!(du, u, p, t)
    # 1. State Recovery (using @views to avoid unrolled ODE heap allocation crashes)
    A_blood = u[1]; A_free = @view u[2:N+1]
    A_surface = @view u[N+2:2N+1]; A_internal = @view u[2N+2:3N+1]
    
    # 2. Fused-Broadcasting (@.) PK Layer math (extracting pure scalar logic out of mapping boundaries)
    k10 = p.k10_pop * (p.GFR / 120.0)
    blood_flow_in = sum(p.f .* p.k_in .* A_blood)
    blood_flow_out = sum(p.k_out .* A_free)
    
    @. p.M_bound = ((A_surface + A_internal) / p.SA_MBq_nmol) * exp(p.lambda_phys * t)
    @. p.binding_flux = p.k3 * A_free * max(0.0, 1.0 - (p.M_bound / p.B_MAX))
    @. p.unbinding_flux = p.k4 * A_surface
    @. p.internalize_flux = p.k5 * A_surface
    
    du[1] = -(k10 + sum(p.f .* p.k_in) + p.lambda_phys) * A_blood + blood_flow_out
    @. du[2:N+1]  = blood_flow_in - (p.k_out * A_free) - (p.lambda_phys * A_free) - p.binding_flux + p.unbinding_flux
    @. du[N+2:2N+1] = p.binding_flux - p.unbinding_flux - p.internalize_flux - (p.lambda_phys * A_surface)
    @. du[2N+2:3N+1] = p.internalize_flux - (p.lambda_phys * A_internal)
    
    # 3. Time-Integrated Activity (TIA) Decoupling
    # TIA includes a vascular fraction so blood-pool irradiates spatial anatomy dimensionally
    @. du[3N+2:4N+1] = A_free + A_surface + A_internal + (p.vol_map * p.v_b * A_blood / p.V_blood)
end

# POST-INTEGRATION SPATIAL DOSE DISCOVERY
final_TIA = sol.u[end].TIA

# Dual Kernel Convolution protecting extensive gamma tracking and unclipped beta bounds
base_energy_beta = compute_base_convolution(final_TIA, p.E_beta, p.K_beta_lung)
base_energy_gamma = compute_base_convolution(final_TIA, p.E_gamma, p.K_gamma_water)
base_energy_total = base_energy_beta .+ base_energy_gamma

# Generate explicit vector flux tracking directionality for neural mapping fields
vector_flux_ray = compute_directional_flux(base_energy_total)

# The NN receives base_energy so it's not blind to cross-fire, out-putting a scalar multiplier
nn_absorbed_fraction_offset = compute_neural_corrector(base_energy_total, p.rho, grad_rho, vector_flux_ray, p.p_nn, p.st_nn)

# Multiply the base_energy multiplicatively by the network offset modifier, then normalize by mass
D_phys = (base_energy_total .* exp.(nn_absorbed_fraction_offset)) ./ mass_map 
\end{lstlisting}

\section{Expert Inquiry and Validated Constraints}

Domain expert review of these formulations has structurally confirmed several native system boundaries natively protecting optimization matrices:
\begin{enumerate}
    \item \textbf{Auger Transmutation Explosions:} While $\sim 15\%$ of $^{177}$Lu decays trigger local internal conversion Auger cascades that fundamentally shred targeting DOTA tags, the physical mass disparity between cold/hot peptides implies evaluating solely a $\lambda_{phys}$ subtraction effectively governs specific activity.
    \item \textbf{Iodine Contrast Shielding Hallucinations:} High-$Z$ Intravenous Contrast ($HU > 300$) scales wildly distinctly via photoelectric phenomena ($Z^3$) compared to pure beta-particle electron densities ($Z/A$). Evaluating iodine maps guarantees mathematical $\beta^-$ attenuation hallucinations, commanding strict Native Non-contrast CT initialization bounds.
    \item \textbf{Directional Vector Embeddings:} Because scaling $\mathcal{N}_\theta$ completely without explicit continuous Vector Directional geometry fields permanently blinds convolutions toward asymmetrical backscatter resolutions, generating discrete Unscattered Ray-Tracing proxy fluxes prior to Deep Learning arrays formally overcomes native radiation symmetries!
\end{enumerate}

\subsection{Questions to Escalate}
\begin{itemize}
    \item \textbf{To Pharmacodynamic Cellular Biologists:} "Historically, our compartmental system forces simple unbinding competitive states ($k_4$) looping into the free interstitial grid boundaries. However, PSMA-ligand configurations map violently toward irreversibly internalized clathrin-coated Endocytic compartments dynamically. Must the macro $A_{bound}$ vector completely separate into surface equilibrium parameters ($A_{surface}$) contrasting rigid integrated sinks ($A_{internal}$) structurally tracking deeply shielded activities?"
    \item \textbf{To Radiobiological Mathematicians:} "Evaluating lethal scaling algorithms across Low Dose Rate trajectories involves severe uncertainties over dynamic repopulation variables. Because we definitively bound individual voxel $T_{treat}(\mathbf{r})$ metrics strictly to lethal duration drop points, should stochastic $G_2$ biological limits and direct Poisson-modeled structural DNA Double Strand Break evaluations overwrite the global $\alpha / \beta$ scalar matrix limits?"
\end{itemize}

\section{Conclusion}
By integrating all variables from `table.txt` including calibration ($CF$), resolution ($RC$), and anatomical mass $m(\mathbf{r})$, this SciML framework provides a truly universal digital twin for precision $^{177}$Lu-PSMA therapy.

\bibliographystyle{unsrt}
\bibliography{bibl}

\end{document}
